% !PRINTED 73
\weberpage{73}

Die Formel (2), auf dies Funktionensystem angewandt, ergibt
die folgende, häufig benutzte Tabelle:

\weq{8}{
\begin{aligned}
\vartheta_{00}(v) &= \vartheta_{01}\!\left(v+\tfrac{1}{2}\right) = \varepsilon\,\vartheta_{10}\!\left(v+\tfrac{\omega}{2}\right) = \varepsilon\,\vartheta_{11}\!\left(v+\tfrac{1+\omega}{2}\right),\\
\vartheta_{01}(v) &= \vartheta_{00}\!\left(v+\tfrac{1}{2}\right) = -i\varepsilon\,\vartheta_{11}\!\left(v+\tfrac{\omega}{2}\right) = i\varepsilon\,\vartheta_{10}\!\left(v+\tfrac{1+\omega}{2}\right),\\
\vartheta_{10}(v) &= \vartheta_{11}\!\left(v+\tfrac{1}{2}\right) = \varepsilon\,\vartheta_{00}\!\left(v+\tfrac{\omega}{2}\right) = \varepsilon\,\vartheta_{01}\!\left(v+\tfrac{1+\omega}{2}\right),\\
\vartheta_{11}(v) &= -\vartheta_{10}\!\left(v+\tfrac{1}{2}\right) = -i\varepsilon\,\vartheta_{01}\!\left(v+\tfrac{\omega}{2}\right) = -i\varepsilon\,\vartheta_{11}\!\left(v+\tfrac{1+\omega}{2}\right),
\end{aligned}
}

worin zur Abkürzung

\[
\varepsilon = e^{\frac{\pi i \omega}{4} + \pi i v}
\]

gesetzt ist.

Nach (3), (4) ist

\weq{9}{
\begin{aligned}
\vartheta_{00}(-v) &= \vartheta_{00}(v), & \vartheta_{01}(-v) &= \vartheta_{01}(v),\\
\vartheta_{10}(-v) &= \vartheta_{10}(v), & \vartheta_{11}(-v) &= -\vartheta_{11}(v),
\end{aligned}
}

d. h. \gesperrt{es sind $\vartheta_{00}(v)$, $\vartheta_{01}(v)$, $\vartheta_{10}(v)$ gerade Funktionen,
$\vartheta_{11}(v)$ ist eine ungerade Funktion.}

Nach § 17 (12) verschwinden die vier Funktionen (7), bzw.
für die folgenden Werte des Arguments

\[
\frac{1+\omega}{2},\quad \frac{\omega}{2},\quad \frac{1}{2},\quad 0,
\]

also für gewisse halbe Perioden. Von Wichtigkeit sind die Werte,
welche die sämtlichen Funktionen (7) für diese Argumentwerte annehmen, und diese lassen sich mittels der Tabelle (8) auf die drei

\[
\vartheta_{00}(0) = \vartheta_{00},\quad \vartheta_{01}(0) = \vartheta_{01},\quad \vartheta_{10}(0) = \vartheta_{10}
\]

zurückführen. Man erhält so aus (8) das folgende System von
Formeln:

\weq{10}{
\begin{aligned}
\vartheta_{00} &= \vartheta_{01}\!\left(\tfrac{1}{2}\right) = \varepsilon_0\,\vartheta_{10}\!\left(\tfrac{\omega}{2}\right) = \varepsilon_0\,\vartheta_{11}\!\left(\tfrac{1+\omega}{2}\right),\\
\vartheta_{01} &= \vartheta_{00}\!\left(\tfrac{1}{2}\right) = -i\varepsilon_0\,\vartheta_{11}\!\left(\tfrac{\omega}{2}\right) = i\varepsilon_0\,\vartheta_{10}\!\left(\tfrac{1+\omega}{2}\right),\\
\vartheta_{10} &= \vartheta_{11}\!\left(\tfrac{1}{2}\right) = \varepsilon_0\,\vartheta_{00}\!\left(\tfrac{\omega}{2}\right) = \varepsilon_0\,\vartheta_{01}\!\left(\tfrac{1+\omega}{2}\right),
\end{aligned}
}
% [sic] The source repeats theta_11 at the far right of row four in (8).
%CONT>>

